% Appendix (single column, after references). Four supporting analyses, all numbers traced to
% notes/experiments.md and data/build_winrate_buckets.py. Same prose register as the body.

\section{The empirical-rate target}
\label{app:teacher}

The reward target $\hat p(x)$ is a state-conditioned empirical win rate estimated from
training-season outcomes. Each play is placed in a bucket defined by three game-level features: the
score margin of the team in possession, in fourteen bins with edges at $\pm1$, $\pm4$, $\pm7$,
$\pm10$, $\pm14$, and $\pm21$ points; the time remaining, in seven bins with edges at $2$, $5$,
$10$, $15$, $30$, and $45$ minutes; and the public pregame point spread, in nine bins with edges at
$\pm0.5$, $\pm3$, $\pm7$, and $\pm10$ points. A bucket's raw value is the fraction of its training
plays whose possession team won. Sparse buckets are shrunk toward their coarser parents by a
hierarchical empirical-Bayes rule: writing $n$ and $w$ for the count and wins in a bucket and
$\hat p_{\text{parent}}$ for its parent estimate, the bucket value is
$(w + M\,\hat p_{\text{parent}})/(n + M)$ with pseudocount $M=25$, applied from the global rate
downward. The estimate for an evaluation play is read from the training-season table alone, so no
test outcome enters its own target.

Finer state features do not raise the target. A fine table that adds field position and down under
the same hierarchical backoff matches the coarse table on game-level calibration: on the 2023 split
the coarse target scores $0.153$ Brier and $0.012$ ECE against the fine table's $0.153$ and $0.010$,
both close to the market's $0.150$ and $0.019$ (\cref{tab:granularity}). Field position and down
shape the current drive rather than the game outcome, so the reward uses the coarse
score-by-time-by-spread table.

\begin{table}[h]
\centering
\small
\begin{tabular}{lcc}
\toprule
Target (eval 2023) & Brier & ECE \\
\midrule
Coarse: score $\times$ time $\times$ spread & 0.1532 & 0.0117 \\
Fine: $+$ field position $+$ down & 0.1534 & 0.0099 \\
Betting market & 0.1495 & 0.0193 \\
\bottomrule
\end{tabular}
\caption{Teacher-granularity ablation on held-out 2023. Adding field position and down to the
coarse buckets does not improve game-level calibration, so the reward target stays coarse.}
\label{tab:granularity}
\end{table}

\section{Training and hyperparameter selection}
\label{app:hparams}

Every checkpoint is selected by Brier score on a fixed random sample of $128$ held-out 2023 states,
scored greedily at each $50$-step save, and the two models are tuned separately. The numbers below
are this in-training selection metric, with one setting varied at a time around a reference
configuration; the reported test numbers use the full splits.

For the direct model the learning rate matters most. Raising it from $5\times10^{-6}$ to
$1\times10^{-5}$ lowers the in-training Brier from $0.160$ to $0.154$ at equal calibration error, and
$2\times10^{-5}$ sharpens further at a higher calibration error; on the larger batch this higher rate
is the optimum. Effective batch size, set by gradient accumulation, improves Brier up to four games per
step and then saturates. Three negatives are decisive (\cref{tab:direct-sweep}): turning off reward
scaling collapses resolution, blending the realized outcome back into the target at $\lambda=0.5$
reintroduces its variance, and lowering the KL coefficient below $0.01$ worsens the calibration
error. Temperature is neutral between $0.7$ and $1.0$, and the selected configuration varies by
$0.004$ Brier across three seeds. The reported direct model uses a learning rate of $2\times10^{-5}$,
gradient accumulation over sixteen micro-batches, a KL coefficient of $0.01$, and reward scaling on.

\begin{table}[h]
\centering
\small
\begin{tabular}{llcc}
\toprule
Lever & Setting & Brier & ECE \\
\midrule
Reference & lr $1{\times}10^{-5}$, 2 games/step, rate target & 0.154 & 0.050 \\
\midrule
\multirow{2}{*}{Learning rate} & $5\times10^{-6}$ & 0.160 & 0.052 \\
                               & $2\times10^{-5}$ & 0.152 & 0.079 \\
\midrule
\multirow{2}{*}{Effective batch} & 4 games/step (grad.\ accum.\ 16) & 0.149 & 0.099 \\
                                 & 8 games/step (grad.\ accum.\ 32) & 0.152 & 0.100 \\
\midrule
\multirow{2}{*}{Reward target} & realized outcome $y$ & 0.166 & 0.065 \\
                               & blend $\tfrac12 y + \tfrac12\hat p$ & 0.181 & 0.121 \\
\midrule
KL coefficient & $0.003$ & 0.157 & 0.100 \\
Reward scaling & off & 0.178 & 0.064 \\
Temperature & $1.0$ & 0.147 & 0.090 \\
\bottomrule
\end{tabular}
\caption{Direct-model ablation, in-training selection metric (greedy, fixed $n=128$ held-out 2023,
best-Brier checkpoint). Each row changes one setting from the reference. The learning rate and the
effective batch size trade sharpness against calibration; reward scaling, the rate target, and a
non-zero KL coefficient are decisive negatives when removed.}
\label{tab:direct-sweep}
\end{table}

The masked model differs in two settings. The KL anchor must be removed: with the gradient
concentrated on the few answer tokens, a coefficient of $0.01$ makes the $k_3$ KL estimator overflow
and the loss diverge, and $1\times10^{-3}$ degrades calibration over training, while $\beta=0$ is
stable because the answer-span mask already protects the reasoning the reference model would
otherwise anchor. The learning rate follows an inverted-U with a higher optimum than for the direct
model (\cref{tab:masked-sweep}): held-out Brier improves from $5\times10^{-6}$ to a peak at
$3\times10^{-5}$ and collapses at $4\times10^{-5}$, where the unanchored update drives probabilities
to extremes. The reported masked model uses a learning rate of $3\times10^{-5}$ and $\beta=0$.

\begin{table}[h]
\centering
\small
\begin{tabular}{llc}
\toprule
Lever & Setting & Brier \\
\midrule
\multirow{3}{*}{KL coefficient} & $\beta=0$ & 0.174 \\
                                & $\beta=1\times10^{-3}$ & 0.190 \\
                                & $\beta=0.01$ & diverges \\
\midrule
\multirow{5}{*}{Learning rate ($\beta=0$)} & $5\times10^{-6}$ & 0.1745 \\
                                           & $1\times10^{-5}$ & 0.1740 \\
                                           & $2\times10^{-5}$ & 0.1717 \\
                                           & $3\times10^{-5}$ & 0.1708 \\
                                           & $4\times10^{-5}$ & over-sharpens \\
\bottomrule
\end{tabular}
\caption{Masked-model ablation, held-out 2024 (n${=}1000$ subset, best checkpoint). The KL anchor
must be removed, and the learning rate follows an inverted-U with an optimum at $3\times10^{-5}$,
above which the unanchored, concentrated gradient drives the maximum-bin error past $0.5$.}
\label{tab:masked-sweep}
\end{table}

\section{Blinded-judge protocol for reasoning quality}
\label{app:judge}

To measure reasoning quality we collect, for each model, its chain of thought on $250$ held-out 2023
plays stratified across game states, with the ground-truth game facts taken from each play's
features. A separate model judges every completion blind to its source and decides, with a quote
from the completion supporting each judgment, whether the chain of thought misreads any of five game
facts (possession, score, spread, clock, and the pregame favorite) and whether the stated
probability follows from the reasoning. We audited the judge by reading every flagged case; each
flag carries a specific checkable reason, and the judge flags the masked model alongside the others.

State-reading errors are low and flat across the base model, the two full-completion reinforcement
runs, and the masked model. Possession and score are misread in about one percent of completions,
the pregame favorite in two to seven percent, the spread in three to six percent, and the clock in
seven to ten percent, the last a base-model weakness that no training variant changes. Reinforcement
learning does not corrupt the reading of the game.

The models separate instead on whether the stated probability follows from the reasoning
(\cref{tab:faithfulness}). The base model is inconsistent on $22.4\%$ of completions, and
full-completion reinforcement learning does not improve it; the failures are broken probability
arithmetic, drive-level and game-level probabilities conflated, and a final number the text does not
support. The masked model is inconsistent on $4.4\%$. The masked prompt format accounts for much of
the reduction, since a base model read under the masked prompt is inconsistent on $6.8\%$, and
masked training lowers it further. At $250$ plays the format and training intervals overlap on this
checkpoint, while a separate control on the reported checkpoint isolates a significant training
effect.

\begin{table}[h]
\centering
\small
\begin{tabular}{lc}
\toprule
Model & Inconsistent reasoning [95\% CI] \\
\midrule
Base (chain of thought) & 22.4\% \ci{17.6}{27.6} \\
Full-completion RL (run A) & 18.8\% \ci{14.4}{23.6} \\
Full-completion RL (run B) & 14.0\% \ci{10.0}{18.4} \\
Base, masked prompt (no training) & 6.8\% \ci{4.0}{9.6} \\
Masked-CoT RLVR (ours) & 4.4\% \ci{2.0}{7.2} \\
\bottomrule
\end{tabular}
\caption{Rate at which the stated probability does not follow from the chain of thought, blinded
judge over $250$ held-out 2023 plays. Full-completion reinforcement learning does not improve on the
base rate; the masked prompt format reduces it sharply and masked training reduces it further.}
\label{tab:faithfulness}
\end{table}

\section{Data-integrity audit}
\label{app:integrity}

The prompt exposes one market-derived quantity, the public pregame point spread, and never the live
in-game win probability that serves as the evaluation reference. We scanned all $40{,}246$ training
prompts and the 2023 and 2024 evaluation prompts for the market win probability and found no
occurrence. The empirical-rate target is built only from realized outcomes and the pregame spread,
and the target for an evaluation play is computed from the training-season table alone, so neither
the market probability nor any test outcome enters training. The splits are disjoint by season,
training on 2015 through 2022, selecting on 2023, and testing on 2024, so no game appears in more
than one split. Every reported number is recomputed from the per-play predictions by a single
deterministic script, so the tables and figures reproduce from the released predictions without a
GPU.
